\scanpage{314}
\mspage{13}
donc
\[
\widehat{\boldsymbol\Delta}^{\,o}\xrightarrow{F_o}
B_o=(\widehat{\boldsymbol\Delta}_{ab})^o
\simeq\underline{\operatorname{Hom}}(\boldsymbol\Delta^o,\mathrm{Ab})^o
\simeq\underline{\operatorname{Hom}}(\boldsymbol\Delta,\mathrm{Ab}^o)
\]
\margin{[$\simeq(\mathrm{Ab}^o)^*\inserted{\simeq(\mathrm{Ab}_*)^o}$ dans le cas qui nous intéresse]}
commutant aux $\varprojlim$
d'où le foncteur
\[
\underline{\operatorname{Hom}}_!(\deleted{\unsure{\widehat{\boldsymbol\Delta}}}B_o,B)
\longrightarrow\underline{\operatorname{Hom}}_!(\widehat{\boldsymbol\Delta}^{\,o},B)
\]
\deleted{et on définit ce foncteur}
dont on veut montrer que c'est une équiv. Posons
$B=A^o$, $B_o=(A_o)^o$, \unsure{on} $F_o=(G_o)^o$, donc
\[
G_o:\widehat{\boldsymbol\Delta}\longrightarrow\widehat{\boldsymbol\Delta}_{ab},\qquad G_o(X)=\mathbb Z(X),
\]
on doit prouver que
\[
\underline{\operatorname{Hom}}^{\deleted{\unsure{!}}}_!(\widehat{\boldsymbol\Delta}_{ab},A)
\longrightarrow\underline{\operatorname{Hom}}^{\deleted{\unsure{!}}}_!(\widehat{\boldsymbol\Delta},A)
\qquad\unsure{\varphi|\psi}\longmapsto\psi\circ G_o
\]
est une équiv., où ${}^{\deleted{\unsure{!}}}_!$ désigne les foncteurs
qui commutent aux $\varinjlim$. Or
le deuxième mb est $\simeq\underline{\operatorname{Hom}}(\boldsymbol\Delta,A)$
(par restriction à $\boldsymbol\Delta\hookrightarrow\widehat{\boldsymbol\Delta}$), \deleted{\unsure{donc}} le
premier s'écrit ? \deleted{\unsure{on}} Prenons $\varphi\in\underline{\operatorname{Hom}}_!(\widehat{\boldsymbol\Delta}_{ab},A)$
et notons que pour tout $M\in A$, \deleted{\unsure{l'objet}} le
foncteur $X\unsure{\vdash|\mapsto}\operatorname{Hom}(\varphi X,M)$ transforme $\varinjlim$
en $\varprojlim$ donc est représentable
\[
\operatorname{Hom}(\varphi X,M)\simeq\operatorname{Hom}(X,\psi M)
\]
i.e. l'hyp. de commutation de $\varphi$ aux $\varinjlim$
s'exprime par le fait qu'il admet un
adjoint $\psi$ --- la condition à mettre sur \deleted{\unsure{\varphi}}
$\psi:A\longrightarrow\widehat{\boldsymbol\Delta}_{ab}$ pour qu'il provienne de $\varphi$
étant que pour tt $X$ fixé dans $\widehat{\boldsymbol\Delta}_{ab}$, $M\longrightarrow\operatorname{Hom}(X,\psi M)$
soit repr. Ecrivant $\widehat{\boldsymbol\Delta}_{ab}\simeq\underline{\operatorname{Hom}}(\boldsymbol\Delta^o,(\mathrm{Ab}))$
on trouve donc interprétation de la
cat. $\underline{\operatorname{Hom}}_!(\widehat{\boldsymbol\Delta}_{ab},A)$ des $\varphi$ comme ss-cat. pleine
de $\underline{\operatorname{Hom}}(A\times\boldsymbol\Delta^o,\mathrm{Ab})$,
\margin{$\simeq\underline{\operatorname{Hom}}(\boldsymbol\Delta^o,\underbrace{\underline{\operatorname{Hom}}(A,\mathrm{Ab})}_{\supset A^o\ \unsure{etc}})$}
foncteurs $\lambda$, satisfaisant aux conditions de
repr. \inserted{ci-dessus} en $A$. Or si $x\in\boldsymbol\Delta$, le foncteur \deleted{$M\longmapsto$}
\deleted{$\lambda(\unsure{M},\unsure{A})$}\inserted{$\lambda(x):A\longrightarrow$}$\in\mathrm{Ab}$ n'est autre que
\[
M\unsure{\vdash|\mapsto}\operatorname{Hom}_{ab}(\mathbb Z(x),\psi(M))
\]
(car
\[
\lambda(\unsure{x})(M)\overset{\mathrm{def}}=\psi(M)(x)
=\operatorname{Hom}_{\widehat{\boldsymbol\Delta}}(x,\psi(M))
=\operatorname{Hom}_{\widehat{\boldsymbol\Delta}_{ab}}(\mathbb Z(x),\psi M)
\]
donc il doit être représentable, \deleted{\unsure{donc}} (faire $X
=\mathbb Z(x)$), donc $\lambda(x)\in A$, donc on est
dans $\underline{\operatorname{Hom}}(\boldsymbol\Delta^o,\deleted{\mathrm{Ab}}A^o)\simeq\underline{\operatorname{Hom}}(\boldsymbol\Delta,A)$,
on a pratiquement terminé --- ]

Quel est \inserted{donc} le foncteur commutant aux
$\varinjlim$ \unsure{purs.} universel de $\widehat{\boldsymbol\Delta}$ dans cat.
additive et avec $\varinjlim$ ? C'est
$X\longrightarrow\mathbb Z(X)=\mathbb Z^{(X)}$ de $\widehat{\boldsymbol\Delta}$ dans $\widehat{\boldsymbol\Delta}_{ab}$
(i.e. le foncteur ``\uline{abélianisation}''), \unsure{donc}
doit se maintenir pour un topos quelconque.
Mais ici \inserted{(i.e. $\boldsymbol\Delta=$ cat. des $\Delta^n$)} cela signifie

\uline{Corollaire} (\unsure{grâce} à Dold-Puppe\deleted{\unsure{pour}}) : \deleted{$\unsure{2^\circ}$}\inserted{$a)$} L'objet
cosimplicial \deleted{$n\longmapsto\mathbb Z_\cdot$} $n\unsure{\vdash|\mapsto}\mathbb C_*({\Delta^n}_*)$
$=\mathbb Z({\Delta^n}_*)=\mathbb Z^{{\Delta^n}_*}$ de $\widehat{\boldsymbol\Delta}_{ab}\simeq\mathrm{Ab}_*$
est \uline{universel} pour les objets cosimpliciaux
dans des cat. additives avec
$\varinjlim$ quelconq. (pour des foncteurs entre
telles cat. qui commutent aux
$\varinjlim$). Donc l'\deleted{objet}\inserted{complexe de chaînes correspondant}
\margin{315}
